\chapter{From a signed history to a capacity institution}\label{ch:institutions}

\section{A record is not yet a permission to spend}
The physical account gives a signed value to a complete transition. A declared attribution can place parts of a joint value beside actors. A further institutional rule might use those records to alter future decision capacity. These are successive choices, not synonyms.

A signed historical account can record harmful and beneficial contributions without requiring its balance to stay nonnegative. A spendable capacity balance can instead impose a nonnegativity rule, a lending rule or another constraint. The potential theorem selects none of them. This distinction is essential if an EBU institution is to support people whose necessary actions sometimes have negative physical value under the current field.

\section{One conditional balance identity}
Suppose a separately declared policy maintains accounts $B_i$ and updates them so that $\sum_i\Delta B_i=E_t$ for the complete actor transition at event $t$. Let $A_{{\mathrm{ext}},t}$ accumulate external potential changes, using
$D_{{\mathrm{ext}},t}=V(x^{\mathrm{b}}_t)-V(x_t)$ before the actor's transition $x^{\mathrm{b}}_t\to x_{t+1}$. Under one fixed field,
\begin{equation}\label{eq:capacity-invariant}
 V(x_t)+\sum_iB_{i,t}-A_{{\mathrm{ext}},t}=K
\end{equation}
is invariant when the updates and initial value $K$ are defined consistently.

\begin{proof}
The change in the left side is
$V(x_{t+1})-V(x_t)+E_t-D_{{\mathrm{ext}},t}$.
Substitute $E_t=V(x^{\mathrm{b}}_t)-V(x_{t+1})$ and the definition of $D_{{\mathrm{ext}},t}$. Every term cancels.
\end{proof}

This is a conditional accounting theorem, not adoption of a wallet. External audit is not a pool that the next actor may spend. Different ownership and update rules can satisfy the same aggregate identity while giving people very different opportunities.

\section{Positive group value does not settle owner permission}
Imagine an attribution rule gives two owners changes $(+3,-1)$, totaling a positive group value of two. Under a proposed rule requiring each post-event capacity to remain nonnegative, two initially empty accounts would refuse the group because the second owner cannot cover the debit. If one owner held both lines and same-owner netting were permitted, the net $+2$ could pass.

The physical endpoint and group EBU have not changed. Ownership and netting changed the institutional decision. A policy that pools all owners, lends capacity or protects an essential action would produce another decision. Those alternatives deserve explicit comparison rather than being hidden inside a routine called settlement.

The possible social benefit of separating these layers is practical: one can retain the physical cost of a necessary service while designing access rules that do not punish its recipient. Hiding the negative physical line would make the system less informative; treating it as an automatic prohibition would make an unproved moral choice.

\section{Zero genesis is a narrow observation}
At a strict quadratic reference, every nonzero net displacement has negative EBU. If all balances begin at zero and a policy requires every account change to be nonnegative, their sum cannot equal that negative group value. Such an action is therefore unaffordable under that particular policy.

This does not prove attraction toward the reference. After a disturbance, positive records may create capacity, and the same policy may permit uphill actions later. A finite menu may also leave no affordable restoring action. A stationary physical state need not imply stationary account balances under continuing external drive.

These possibilities explain why a capacity mechanism requires a dynamical study beyond its invariant. Closure constrains the process but does not choose its behavior. A correctly implemented mechanism can cycle, ration, refuse or perform badly on service.

\section{Repeated damage and repair need complete provenance}
Under a fixed field, a complete closed physical cycle has zero aggregate EBU. This can block one simple accounting exploit when every signed transition is recorded. It does not by itself identify who caused damage, prevent an actor from concealing the negative event, or guarantee that gains and losses are assigned to the same owner.

A person who damages a shared asset and another who repairs it may have opposite historical records. Whether compensation or responsibility follows needs evidence and governance. The aggregate theorem cannot infer intent from endpoints. Essential maintenance can also recur because of genuine external wear or demand, so repeated positive restoration is not automatically suspicious.

A useful institution would therefore preserve event identities, causal evidence where available, correction procedures and appeals. These are social and engineering requirements supporting the account, not conclusions secretly contained in the equation.

\section{Atomic commitment protects a joint account}
If an institution authorizes one joint group, its physical and account records must refer to that exact group. Crediting one child and losing another child's debit can manufacture a claim. Executing the physical action twice after a receipt failure can duplicate the consequence. An eventual operating protocol needs durable identity and explicit recovery semantics.

Here ``atomic'' means a coherent all-or-nothing institutional commitment under the declared protocol. It does not mean concealing a chain of physical transformations inside one undifferentiated event. Fuel, electricity, transport and water effects still need traceable boundaries.

No operating protocol is adopted by this discussion. The practical feature is a requirement that an implementer can inspect: a committed entitlement must be linked to the authorized complete event and applied once. If implemented and audited, that feature could reduce both accidental duplication and deliberate overclaiming.

\section{Access changes whether information can help}
A household may be unable to replace an inefficient heater even when the physical benefit is obvious. A remote clinic may depend on a costly route because no alternative exists. Giving those actors better scores does not give them capital, infrastructure or bargaining power.

An EBU institution could use physical evidence to help target public investment, support maintenance or identify shared services. It could also worsen exclusion if adverse records become automatic barriers to care or mobility. The same accurate measurement can enter very different institutions.

The constructive question is therefore how the information changes feasible opportunities. Does it help finance a repair that lowers future burden? Does it reveal an unrepresented service failure? Can affected people challenge the field or the measurement? Are people with unavoidable needs protected? These are concrete criteria for later institutional studies, not benefits established by the balance identity.

\section{A place for learning without retrospective rewriting}
If a field turns out to omit an important consequence, revise it through a visible process. If an allocation produces unfair outcomes, revise that rule while retaining its historical version. The separation of physical value, attribution and capacity makes such changes more intelligible: one can identify which layer failed and what evidence motivated the revision.

That separation is a strength of the proposed architecture. It gives a future society a way to learn without pretending that either a mathematical theorem or a political rule was infallible. The next chapter translates the same discipline into the responsibilities of an implementation.
